\section{Introduction}
\label{sec:intro}
\FigurePath

Politicians often answer the question they would have liked to be asked.
\citet{thomas-etal-2024-never} turned this observation into a two-level taxonomy of response clarity and evasion, and SemEval-2026 Task~6 \citep{thomas-etal-2026-semeval} made it a shared task.
The task was won by pipelines that call a large language model (LLM) several times per item \citep{shi-etal-2026-teleai}, while fine-tuned encoders levelled off near 0.50 macro-F1 on the nine-way subtask.
A leaderboard shows which system is best, but not why the others stop where they do, nor what a good classifier of evasion could be used to find out.

This paper addresses those two questions in five stages (Figure~\ref{fig:path}). Figure~\ref{fig:system} shows the system it builds and how it is used.
At each stage we wrote down a hypothesis, a prediction and a decision rule before running anything, and the result of each stage set the question of the next.

\noindent Our contributions follow the same order. First, controlled evidence that a stronger pretrained backbone helps where structural changes do not. Second, a cascade at 0.650 Subtask-2 macro-F1 with 1.4 LLM calls per item, and a distilled variant at a quarter of the classifier's GPU time. Third, evidence that evasion is relational, and an exploratory LLM audit in which, in the President's role, models answer directly what the Presidents answered directly. Fourth, a measurement of what reduced precision costs the cascade.
Every system comparison is on the 308-item development set, since the test labels were never released.
Many registered hypotheses failed. We report them where they fell, because each failure changed what we did next.
